% ═════════════════════════════════════════════════════════════════════════════
\section{View Abstraction}
\label{sec:views}
% ═════════════════════════════════════════════════════════════════════════════

Monotonic abstraction computes the configurations from which a bad
configuration is reachable, starting from the bad configurations and going
backward. In this section, we present \emph{view
abstraction}~\cite{AHH13,AHH16}, a forward analysis that exploits a
\emph{small model property} of many parameterized systems: the behavior of
large instances is already visible in small ones. A configuration of a large
instance can often be seen as a combination of patterns that already occur in
configurations of small instances, and if a bad pattern exists, it can often be
found with only a few processes. The method therefore inspects only small
instances of the system, and abstracts the configurations it finds into
\emph{views}, which describe the system from the perspective of a fixed number
of processes. The views are then recombined to obtain an over-approximation of
the reachable configurations of all sizes. The number of processes that need to
be inspected is not fixed in advance: the method discovers a \emph{cutoff
point} during the analysis, beyond which the search need not continue. In this
sense, view abstraction is a \emph{dynamic cutoff} method, in contrast to the
static cutoff theorems discussed in Section~\ref{sec:other}. Through the theory
of Section~\ref{sec:wqo}, the method is complete for a large class of
well-structured systems.

\subsection{Model}

We consider parameterized systems with a linear topology as defined in
Section~\ref{sec:prelim}, and we assume that global conditions are checked
atomically. For instance, the rule~$t_4$ of Figure~\ref{fig:mutex} is written
$\textbf{if}\ \forall j < i : \{\Idle\}\ \textbf{then}\ \Wait \to \Crit$. We assume that $\Bad = \upc{B_{\min}}$ is the
upward closure of a finite set $B_{\min}$ of minimal bad configurations with
respect to the subword ordering. We write $\mathcal{R}_k$ for the set of configurations of size
at most~$k$ that are reachable from initial configurations of size at
most~$k$. Since the transitions do not change the size of a configuration,
$\mathcal{R}_k$ is finite and can be computed by explicit-state exploration.

\subsection{Views}

For a natural number~$k$, a \emph{view} of size~$k$ of a configuration is a
subword of the configuration of length~$k$. A view retains the states of $k$
processes, in the order in which they occur in the configuration, and ignores
the other processes. We write $\mathcal{V}_k$ for the set of nonempty words of
length at most~$k$ over~$Q$, which are the possible views of size up to~$k$. The
\emph{abstraction function} $\alpha_k$ maps a configuration~$c$ to the set of
its views of size up to~$k$, that is, to the set of words $v \in
\mathcal{V}_k$ with $v \preceq c$. For a set~$X$ of configurations, $\alpha_k(X)$
is the union of the sets $\alpha_k(c)$ with $c \in X$. The \emph{concretization
function} $\gamma_k$ maps a set~$V$ of views to the set of all configurations
that can be built from views in~$V$:
\[
\gamma_k(V) = \{ c \mid \alpha_k(c) \subseteq V \} .
\]
The two functions form a Galois connection, so $X \subseteq
\gamma_k(\alpha_k(X))$ for every set~$X$ of configurations.

The concretization is an over-approximation. Consider the set $X =
\{\Idle\Req\Wait,\, \Req\Wait\Crit,\, \Idle\Wait\Crit\}$. Its views of size up
to two are the four single states together with $\Idle\Req$, $\Idle\Wait$,
$\Req\Wait$, $\Req\Crit$, $\Wait\Crit$, and $\Idle\Crit$. The concretization of
these views contains the three configurations of~$X$, but also, for instance,
the configuration $\Idle\Req\Crit$, all of whose views of size two belong to
the set, and the configuration $\Idle\Req\Wait\Crit$ of size four. The views
work together to represent a set of configurations: a configuration belongs to
the concretization if and only if all its views of size up to~$k$ belong to
the set. Larger views impose more
constraints, and the precision increases with~$k$: for every set~$X$ of
configurations, we have $\gamma_{k+1}(\alpha_{k+1}(X)) \subseteq
\gamma_k(\alpha_k(X))$.

The \emph{abstract post-image} of a set~$V$ of views is defined as usual by
$\mathit{Apost}_k(V) = \alpha_k(\Post(\gamma_k(V)))$. If a set~$V$ of views
contains the views of the initial configurations and satisfies
$\mathit{Apost}_k(V) \subseteq V$, then $\gamma_k(V)$ contains all reachable
configurations, that is, $\mathcal{R} \subseteq \gamma_k(V)$.

\subsection{The Procedure}

\begin{algorithm}[t]
\caption{Verification by view abstraction}
\label{alg:views}
\begin{algorithmic}[1]
\FOR{$k = 1, 2, 3, \ldots$}
  \IF{$\mathcal{R}_k \cap \Bad \neq \emptyset$}
    \RETURN unsafe
  \ENDIF
  \STATE $V \leftarrow$ the least set of views that contains $\alpha_k(\Init)$ and satisfies $\mathit{Apost}_k(V) \subseteq V$
  \IF{$\gamma_k(V) \cap \Bad = \emptyset$}
    \RETURN safe
  \ENDIF
\ENDFOR
\end{algorithmic}
\end{algorithm}

The verification procedure (Procedure~\ref{alg:views}) consists of two nested
loops. The outer loop increases the precision~$k$. For a given~$k$, the
procedure first computes the reachable configurations of size up to~$k$
exactly. If one of them is bad, the system is unsafe, and the run that reaches
it is a counterexample. Otherwise, the procedure computes a fixpoint of the
abstract post-image, starting from the views of the initial configurations. If
the concretization of the fixpoint contains no bad configuration, the system is
safe for all sizes. Otherwise, the result is inconclusive, since the bad
configuration may be an artifact of the abstraction, and the procedure
increases~$k$. The value of~$k$ for which the procedure returns safe is a
dynamic cutoff.

The fixpoint computation cannot be performed directly, since $\gamma_k(V)$ is
infinite in general. The key observation is that it suffices to consider the
configurations of size at most $k+1$ in $\gamma_k(V)$:
\[
\alpha_k(\Post(\gamma_k(V))) \cup V
= \alpha_k(\Post(\gamma_k^{k+1}(V))) \cup V ,
\]
where $\gamma_k^{k+1}(V)$ is the finite set of configurations of size at most
$k+1$ in $\gamma_k(V)$. The reason is that a transition changes only the state
of the active process, so a view of the new configuration that is not already
in~$V$ must contain the active process. Consider the subword of the old
configuration that consists of the processes of such a view. This subword
belongs to $\gamma_k(V)$, since its views are views of the old configuration.
For a local or a universal rule, the transition is still enabled in this
subword, since removing processes can only make a universal condition easier
to satisfy. For an existential rule, we also keep one witness of the old
configuration, which gives a subword of the old configuration of size at most
$k+1$. It belongs to $\gamma_k(V)$ for the same reason, and the transition is
enabled in it. In both cases, the view is produced from a
configuration of size at most $k+1$ in $\gamma_k(V)$. The test $\gamma_k(V)
\cap \Bad = \emptyset$ is also finite: it amounts to checking that no configuration $b \in B_{\min}$ has all
its views of size up to~$k$ in~$V$. In practice, the fixpoint computation is
accelerated by starting from the views of $\mathcal{R}_{k+1}$, which is
obtained by explicit-state exploration of the instances with at most $k+1$
processes, rather than from the views of~$\Init$. In most experiments, this
initial set is already a fixpoint or close to one~\cite{AHH16}.

\paragraph{Example.}
We apply the procedure to the protocol of Figure~\ref{fig:mutex}. For $k = 1$,
the reachable configurations of size one contain no bad configuration, and the
fixpoint consists of the four views $\Idle$, $\Req$, $\Wait$, and $\Crit$. Its
concretization contains all configurations, including $\Crit\Crit$, so the
result is inconclusive. For $k = 2$, the reachable configurations of size at
most two contain no bad configuration. The views of the reachable
configurations of size at most three are all nonempty words of length at most
two except $\Crit\Crit$, $\Req\Crit$, and $\Wait\Crit$. Starting from these
views, as described above, the abstract post-image, computed from the
configurations of size at most three in their concretization, produces no new
view, so the fixpoint is reached immediately. Since $\Crit\Crit$ is not among the views, the
concretization contains no bad configuration, and the protocol is safe for all
sizes, with dynamic cutoff~$2$. The fixpoint describes exactly the
configurations that contain none of $\Crit\Crit$, $\Req\Crit$, and $\Wait\Crit$
as a subword. These are the complement of the set computed by monotonic
abstraction in Section~\ref{sec:monotonic}, which was generated by the same
three configurations. The two methods arrive at the same invariant, one going
forward from the initial configurations and the other going backward from the
bad ones.

\subsection{Termination and Completeness}

The procedure is sound: when it returns safe, the system is safe, and when it
returns unsafe, the counterexample is real. If the system is unsafe, the
procedure terminates, since a bad configuration that is reachable from an
initial configuration of size~$n$ belongs to~$\mathcal{R}_k$ for every
$k \geq n$. For safe systems, termination cannot be guaranteed in general,
since the problem is undecidable. However, the theory of well-quasi-orderings
of Section~\ref{sec:wqo} identifies a large class of systems for which the
procedure always terminates. We show this in three steps.

We first establish when the reachable set~$\mathcal{R}$ can be represented
exactly by views. For a set~$X$ of configurations, we define its
\emph{downward closure} ${\downarrow}X$ as the set of configurations~$c$
with $c \preceq c'$ for some $c' \in X$, and we say that $X$ is
\emph{downward-closed} if $X = {\downarrow}X$.
Suppose that $\mathcal{R}$ is downward-closed. Its complement is then
upward-closed, and by Higman's lemma it has a finite generator. Let $K$ be the
maximal size of a configuration in this generator. Then
$\gamma_K(\alpha_K(\mathcal{R})) = \mathcal{R}$: a configuration
outside~$\mathcal{R}$ contains an element of the generator as a subword,
which is a view of size at most~$K$ that does not occur in any reachable
configuration. Thus, $\mathcal{R}$ can be represented exactly by views of
size at most~$K$.

Second, monotonicity guarantees a sufficiently precise representation even
when $\mathcal{R}$ is not downward-closed. If the transition relation is
monotonic with respect to the subword ordering and $\Bad$ is upward-closed,
then ${\downarrow}\mathcal{R}$ is closed under transitions, and a bad
configuration is reachable if and only if ${\downarrow}\mathcal{R}$ contains
one. Thus, if the system is safe, ${\downarrow}\mathcal{R}$ is a
downward-closed inductive invariant that proves safety. By the first step,
applied to ${\downarrow}\mathcal{R}$, it can be represented exactly by views of
size at most~$K$, where $K$ is now the maximal
size of a configuration in the generator of the complement
of~${\downarrow}\mathcal{R}$. This covers systems with local and existential
transitions, and, with the extensions discussed below, Petri nets and
broadcast protocols. The result extends to general well-structured transition
systems, such as lossy channel systems, provided that the configurations can be
measured by a size that has finitely many configurations of each
value~\cite{AHH16}.

Finally, the procedure does indeed compute this representation. In both
cases, the set ${\downarrow}\mathcal{R}$ is closed under transitions, and it
is equal to~$\mathcal{R}$ in the first case. Since
$\gamma_K(\alpha_K({\downarrow}\mathcal{R})) = {\downarrow}\mathcal{R}$, the
set $\alpha_K({\downarrow}\mathcal{R})$ contains the views of the initial
configurations and is closed under the abstract post-image. Since the
procedure computes the least such set, its result~$V$ for $k = K$ is contained
in $\alpha_K({\downarrow}\mathcal{R})$. Conversely, soundness gives
$\mathcal{R} \subseteq \gamma_K(V)$, and hence $\alpha_K(\mathcal{R})
\subseteq V$. Taking the downward closure introduces no new views, so
$\alpha_K(\mathcal{R}) = \alpha_K({\downarrow}\mathcal{R})$, and we conclude
that $V = \alpha_K({\downarrow}\mathcal{R})$. Hence $\gamma_K(V) =
{\downarrow}\mathcal{R}$, which contains no bad configuration since the system
is safe, and the procedure returns safe at the latest when $k = K$.

For these systems, view abstraction provides a forward alternative to
the backward reachability analysis of Section~\ref{sec:wqo}.
The protocol of Figure~\ref{fig:mutex} shows that the method can also
succeed without monotonicity: its universal transitions are not monotonic,
but its reachable set is downward-closed. The generator of its complement
is $\{\Crit\Crit, \Req\Crit, \Wait\Crit\}$, giving $K = 2$.

\subsection{Extensions}

The method extends to other communication primitives and topologies by
adapting the abstraction function, the concretization, and the post-image.
For broadcast transitions, the process that initiates the broadcast plays the
role of the witness, and configurations of size at most $k+1$ still suffice.
A rendezvous transition requires the presence of several processes, and the
configurations used in the abstract post-image must be of size up to $k+m$
when a rendezvous involves $m+1$ processes. Shared variables are modeled by a controller process whose
state is kept in every view. For multisets, views are sub-multisets; for
rings, they are circular subwords; and for trees, they are subtrees of a given
depth~\cite{AHH16}.

Some safe systems have no downward-closed inductive invariant that excludes the bad
configurations, and their analysis with the views above produces spurious
counterexamples for every~$k$. A typical case is
a universal transition that is disabled in all reachable configurations by a
process that a view of size~$k$ ignores. To handle such systems,
\emph{context-sensitive views} record, in addition to the $k$ retained
processes, the sets of states of the ignored processes between them. These
contexts are taken into account when universal conditions are evaluated, which
blocks the spurious transitions~\cite{AHH14,AHH16}. A further refinement of the
views handles global conditions that are not checked atomically, by recording
for each process how far it has progressed in its inspection of the others.
With these extensions, the method verified the full, non-atomic version of
Szymanski's mutual exclusion protocol, which, to the knowledge of the authors
of~\cite{AHH16}, had not been verified automatically before.
In the experiments reported in~\cite{AHH16}, covering cache coherence
protocols, protocols on rings and trees, and Petri nets with inhibitor
arcs, the method terminated almost immediately on most examples,
suggesting that the relevant patterns already occur in small instances.
View abstraction has also been extended to programs with more complex
states, such as the heap shapes of lock-free concurrent data
structures~\cite{AHHJR17}.

